\pdfmapfile{+cm.map}
\pdfmapfile{+cmextra.map}
\pdfmapfile{+symbols.map}
\documentclass[11pt]{article}
\usepackage[margin=0.82in]{geometry}
\usepackage{amsmath,amssymb,amsthm,longtable,booktabs,array,hyperref}
\hypersetup{colorlinks=true,linkcolor=blue,urlcolor=blue,pdftitle={Joint sign guards and Pell descent for smaller consistency equations},pdfauthor={OpenAI-assisted research for Vladimir Reshetnikov}}
\renewcommand{\labelitemi}{$\bullet$}
\newtheorem{theorem}{Theorem}\newtheorem{lemma}{Lemma}
\providecommand{\mathscr}[1]{\mathcal{#1}}
\newcommand{\Z}{\mathbb Z}\newcommand{\N}{\mathbb N}
\newcommand{\Con}{\operatorname{Con}}
\setlength{\emergencystretch}{4em}
\title{Joint sign guards and Pell descent\\for smaller consistency equations}
\author{OpenAI-assisted research for Vladimir Reshetnikov and ProveIt}
\date{24 October 2026 (UTC)}
\begin{document}\maketitle
\begin{abstract}
This iteration combines sign recovery, calibrated native guards, simultaneous
positivity, and carefully ordered elimination of arithmetic coordinates.
The resulting unrestricted-integer family includes a quartic in 37 unknowns,
a sextic in 31, an octic in 25, and polynomials in 23, 22, 21, 20 and 19
unknowns at degrees 12, 16, 20, 32 and 88 respectively. The important gain is not a new universal-machine bound: it is
that joint sign compression and zero-set row reduction allow many more
source coordinates to disappear without introducing rational witnesses.
We give complete fixed coefficient bindings and compact defining equations.

Independently, adjacent-coefficient Pell descent removes a denominator guard
from the eleven-positive-variable construction, lowering degree 325 to 277
and operations 113 to 111. A centered transport gives degree 257 in 114 operations, a further degree/cost
tradeoff. The direct proof-checker route also gains smaller exact local
polynomial graphs, and the inert-prime certificate no longer needs a CRT
index variable or a second variable-power graph. Those compiler-interface
advances are not presented as completed small consistency equations.

All endpoint arguments are ordinary mathematics with explicit inherited
Pell, native compiler and refutation-calibration dependencies. Finite replays
and independent symbolic checks supplement those arguments; they do not
constitute a new end-to-end Lean/Rocq certificate or a proof of consistency.
No global minimality or literature-wide priority is claimed.
\end{abstract}
\tableofcontents
\section{The target and the distinction between an equation and an interface}
We seek a fixed polynomial $P\in\Z[x_1,\ldots,x_n]$ and a stated domain $D$
for which
\[
\mathrm{ZFC}\vdash
\bigl(\neg\exists\boldsymbol{x}\in D^n\ P(\boldsymbol{x})=0\bigr)
\longleftrightarrow\Con(\mathrm{ZFC}).
\]
Here $\N$ includes zero; $\N_{>0}$ and $\Z$ are different domains.
A polynomial register is an abbreviation, not a hidden witness. Fixed
coefficient integers are specified independently of consistency, by complete
finite recipes and authenticated payloads. None of our defining equations
contains an unknown exponent, factorial or binomial coefficient.
The alternative CRT interface does use such functions and is explicitly
excluded from the completed-polynomial rankings.

The constructions are compared along several axes: domain, unknown count,
degree, arithmetic operations, coefficient payload and proof burden.
An unrestricted-integer equation is not dominated by a positive-domain
equation merely because their nominal variable counts agree.
Likewise a reduction in variables at the cost of degree is a tradeoff,
not an unqualified simplification of every representation.
Our principal completed family uses the coherent source34 tuple and two
additional powers of two. The positive Pell family uses compact44 instead.
No individual coefficient is exchanged between these tuples.

The fresh starting point is repository commit
\texttt{2605669a532342f153b13eb6b4a0af4b54f5fa37}.
All eight original wave-ten archives are pinned and retained unchanged.
The ordinary proofs establish the reductions relative to the named inherited
results, not a newly checked derivation in the repository's encoded ZFC
calculus. This distinction does not turn a finite test into a universal
proof, nor does it assume the truth of $\Con(\mathrm{ZFC})$.
\section{Fresh baseline and proof boundary}
The immutable repository checkpoint is \texttt{2605669a5323}; the eight
wave-ten originals are pinned to arrival \texttt{9513bd88cefb}. We checked
all 1,183 release-manifest entries and replayed the eight roots normally and
seven in optimized Python, always in fresh temporary copies. All fifteen
replays passed, and every original remained unchanged. A separate reader
verified topology, liveness and operation counts for 21 literal circuits.
These are finite algebra/data checks, not a new proof-assistant build or an
internal consistency derivation.

\begin{center}
\begin{tabular}{lrrr}
Domain and family&Variables&Degree&Operations\\\hline
Integer, cone or determinant&38&4&not optimized\\
Integer, cone or determinant&32&6&not optimized\\
Integer, cone or determinant&27&8&not optimized\\
Integer, nested positivity&26&16&not optimized\\
Integer, free Lucas&10&1,576,864&481\\
Positive, root transfer&11&325&113\\
Positive, rational Pell descent&12&82&102\\
Positive, centered product&12&110&100\\
Positive, centered SOS&12&164&99\\
Positive, strong congruence&13&28&98\\
Positive, adaptive carrier&13&30&92\\
Positive, retained product&14&24&101\\
Positive, composed nine&9&1,379,730&478\\
Positive, repository H2082&9&1,845,940&443
\end{tabular}
\end{center}
This is a selected domain-preserving comparison, not a global optimality
claim. An operation is one live binary addition, subtraction or
multiplication; fixed numeral leaves and register reuse are free.
The cone and determinant equations have the same sizes but different
polynomials. Their six fixed payloads have respectively 881,696 and
615,883 raw bytes. Neither licenses arbitrary substitution of the smaller
compact54 tuple.

The latest formal checkpoint completes the supplied numeric sealed-ZFC
refutation checker with 83,281 raw and 76,313 optimized program-tree nodes.
Those are source-syntax counts, not polynomial sizes. Native arithmetic now
extracts sparse assignments and the local clause bands, but whole-payload
binding, phase synchronization and global native calibration remain open.
The internal semiring, local wire quartics, sum-of-squares conjunction and
literal tree-to-quartic compilation are checked on all internal naturals.
The full checker transport on nonstandard proof codes and the final
object-language consistency/unsolvability biconditional remain open.

\section{Commuting restrictions: 37 integers at degree four}
The degree-four, degree-six and degree-eight unrestricted-integer endpoints
improve from the inspected incoming sizes $38,32,27$ to
\[
 \boxed{37\text{ variables, degree }4},\qquad
 \boxed{31\text{ variables, degree }6},\qquad
 \boxed{26\text{ variables, degree }8}.
\]
This is a compatibility result: source 70's residue orientation and source 71's
calibrated determinant guard restrict the \emph{same} canonical source-34
witness family, so their savings combine. Neither ingredient alone is new here.
The conclusions use the explicitly inherited unbounded Pell, native compiler
and syntactic-calibration theorems. They are not new kernel-checked consistency
biconditionals, and no global minimality is claimed.

\subsection{The explicit 37-variable quartic}
Use the authentic fixed tuple $(\beta,K_0,M_C,M_F)$, with
$\beta=2^{1064921}-1$ and $M_F$ already shifted by $\beta$.
The only new fixed numerals are
\[
 A_0=2^{1113},\qquad B_0=2^{666104}.
\]
All six integers are regenerated from their recipes and actual payload bytes;
their total minimal unsigned size is $615{,}883$ bytes.
Let
\[
 q=\beta J+1,\quad C=q-F-\alpha,\quad X=3X_3-1,\quad
 H=4a+3,\quad k=\eta+\zeta,
 \quad \operatorname{Pos}(x,y,z)=1+x^2+y^2+z(z+1).
\]
These are polynomial abbreviations, not additional unknowns.
The 37 unrestricted integer unknowns are
\[
\begin{gathered}
 J,F,\alpha,z,f,h,r,v,w,\tau,\eta,\zeta,y,\gamma,o,j,
 N,E,P,c,a,X_3,D,C_2,L,F_2,V,g,I,B_{\rm aux},\Delta,\\
 e_1,e_2,e_3,t_1,t_2,t_3.
\end{gathered}
\]
Here $v,w$ abbreviate the JSON names \texttt{v\_scale,w\_scale},
$e_1,e_2,e_3$ the retained $\eta\zeta$ positivity coordinates,
and $t_1,t_2,t_3$ the new guard coordinates.
Define the following 24 residuals by subtracting each right side from its
left side:
\begin{align*}
 L&=q^2-1,&N&=q(L+1),&X&=4qw,\\
 a-E&=(2v+1)N,&E(X+1)&=Xa,&c&=(a-E)k+\eta,\\
 P&=k+E(c-\eta),&\tau^2&=P^2-Pk+1,&D&=X+ac+\gamma H,\\
 \Delta&=a^2+4a+3,&C_2&=c^2,&D^2-\Delta C_2&=1,\\
 F_2&=f^2,&V+c&=fo,&V-cj+k-1&=hE,\\
 g&=y^2-V^2,&I&=(2r+1)C_2,&B_{\rm aux}&=I^2,\\
 B_{\rm aux}&=\Delta(F_2-1),&y^2-B_{\rm aux}g&=1.
\end{align*}
The remaining four equations are
\begin{align*}
 0={}&k-hE+FN+qL-qN+q^2-1-F(q+L)\\
 &{}-\alpha L-(M_C+qM_F)J,\\
 (K_0+4w)C={}&F+z(q-1),\\
 \eta\zeta={}&\operatorname{Pos}(e_1,e_2,e_3),\\
 (A_0F)(B_0C)-(q-A_0F-B_0C)^2={}&\operatorname{Pos}(t_1,t_2,t_3).
\end{align*}
The polynomial is the sum of these 24 squared residuals. All residuals are
quadratic or lower in the displayed unknowns, and the coefficient of
$\tau^4$ is exactly one. Thus the polynomial has exact degree four after
authentic coefficient specialization. Its complete expanded residual list is
\texttt{determinant37\_quartic\_direct.json}.

\subsection{Arbitrary-signed soundness, before canonicality}
Put $U=A_0F$ and $V_0=B_0C$.
The determinant row gives $UV_0>(q-U-V_0)^2\ge0$.
Consequently $U,V_0$ have the same nonzero sign and
\[
 |q-U-V_0|<\sqrt{UV_0}\le\frac{|U|+|V_0|}{2}.
\]
Thus $q,F,C$ have one sign, and
$|q|>(A_0|F|+B_0|C|)/2\ge|F|+|C|$ forces $\alpha$ to share it.
This restores both old product guards using two fresh Pos triples.
No canonical-witness or positive-radix assumption has been made.

It remains to restore the deleted square witness, not to assume it.
Set $Q=|q|$, $Y=(2v+1)q^3$ and $R$ equal to the authentic packing index.
The retained lift yields
\[
 E=XY,\quad a=(X+1)Y,\quad c=kY+\eta,\quad
 k=R+1+hE,\quad V=of-c=jc-R.
\]
The common-sign packing estimates imply
$Q\ge16$ and $Q(Q^2-1)<R<Q^4-2$.
Since $X=3X_3-1\ne0$ and $X=4qw$,
\[
 |X|\ge4Q,\qquad |Y|\ge Q^3,\qquad |E|\ge4Q^4>3R+1.
\]
Write $B=|a+2|$, $P_0=|2XY^2+1|$, $c_0=|c|$, $K=|k|$.
The first and main Pell norms, valid for either sign of $X$, give
\[
 K=2\psi_{P_0}(n),\qquad c_0=\psi_B(p),\qquad p>n.
\]
Here $P_0>B>R$ and $K(|Y|-1)<c_0<K(|Y|+1)$.
The index congruence is
\[
 \operatorname{sgn}(k)\operatorname{sgn}(X)^{n-1}\,2n
 \equiv R+1\pmod {|E|}.
\]
Its positive branch has $n\ge(R+1)/2$ and its negative branch has $n>R$.
These estimates establish the inherited ordinary strong-rank hypotheses.
The signed auxiliary step-down then gives $p=R$ without choosing any sign.
The negative branch is impossible, and no wrap is possible in the other:
$R=2n-1$ is odd.

For $U_0=|X|\ge64$, $Z_0=|Y|\ge4096$ and $r_0=(R-1)/2$,
Pell growth gives
\[
 \frac{c_0}{K}\ge\frac12
 \left(\frac{(2B-1)^2}{2P_0}\right)^{r_0}
 >\frac{(U_0-3)^{r_0}}2.
\]
The strict inequality follows from the universally positive difference
\[
 (2(U_0-1)Z_0-5)^2-(U_0-3)(4U_0Z_0^2+2)>0.
\]
After $U_0=64+u$, $Z_0=6+z$ it is
$4uz^2+28uz+22u+260z^2+1860z+1703$.
Hence $Z_0>2^R$.
The signed Pell identities modulo the odd $H$ give either
$X2^R=\pm1\pmod H$ or $X=e2^R\pmod H$, $e\in\{\pm1\}$.
The strict bounds $1<|X|2^R<|H|-1$ and
$|X-e2^R|<|H|$ exclude the reciprocal branch and unwrap the direct one.
Therefore $X=e2^R$ with odd $R$.
Only now does $X\equiv2\pmod3$ force $e=1$.
Restoring $u_{\rm old}=2^{(R-1)/2}$ recovers the missing row
$X=2u_{\rm old}^2$ and every original Q41 equation.
The inherited Q41/native/calibration theorem supplies the finite refutation.

\subsection{Why completeness survives both restrictions}
The canonical source-34 witnesses have $q=2^t$ with odd $t$,
$R\equiv3\pmod4$, and $X=2^R$; thus $(X+1)/3$ is integral.
The new determinant holds for this same canonical family, including all
permitted rank-controller dummy changes. In radix $\rho=2^{121}$,
each occupied cell has its fixed top $C$ digit at position 3296 and its
fixed top $F$ digit at position 8791, of coefficient
$c_0=105534389824824194647382298625$.
The carry-free native bounds and literal scale identities
\[
 121\cdot8791+1113=1064921-97,\qquad
 121\cdot3296+666104=1064921-1
\]
imply, uniformly in tableau size and rank tuning,
\[
 q/2\le U<4q/5,\qquad q/2\le V_0<3q/5.
\]
Writing $q=2s$, $U=s+x$, $V_0=s+y$ gives
$0\le x<3s/5$, $0\le y<s/5$, so
\[
 UV_0-(q-U-V_0)^2
 =s^2+s(x+y)-x^2-xy-y^2>\frac{12}{25}s^2>0.
\]
A Pos representation supplies the three new guard coordinates.
Thus the same canonical witnesses satisfy both restrictions simultaneously.
This does not assert that every old zero satisfies the new determinant.

\subsection{The sextic and octic companions}
For projections use the equally valid affine quartic
$J_{\rm old}=3J_3+1$, $w_{\rm old}=3w_3+1$ rather than replacing $X$.
It still has $X\equiv2\pmod3$.
Canonical completeness holds because $\beta\equiv1\pmod3$ gives
$J_{\rm old}\equiv1\pmod3$, while
$w_{\rm old}=2^{R-t-2}\equiv1\pmod3$ has a nonnegative even exponent.
Put $q=\beta(3J_3+1)+1$ and $W=4(3w_3+1)$.

For the sextic set $\tau_{\rm old}=\tau_{\rm tail}+E(c-\eta)$ and restore
\[
 P=k+E(c-\eta),\quad N=q(L+1),\quad X=Wq,\quad
 C_2=c^2,\quad F_2=f^2,\quad g=y^2-V^2.
\]
Six coordinates and six defining rows disappear. The first norm becomes
$\tau_{\rm tail}^2+(2\tau_{\rm tail}-k)E(c-\eta)-1$.
All 18 residuals in 31 variables have degree at most three, and their
sum of squares has coefficient four at $r^2c^4$.

For the octic put $k_0=1-V+jc+hE$ and restore
\[
\begin{gathered}
 L=q^2-1,\quad N=q^3,\quad X=Wq,\quad C_2=c^2,\quad F_2=f^2,\\
 g=y^2-V^2,\quad B_{\rm aux}=I^2,\quad
 D=Wq+ac+\gamma(4a+3),\quad\Delta=a^2+4a+3,\\
 \zeta=k_0-\eta,\qquad P=k_0+E(c-\eta).
\end{gathered}
\]
Eleven coordinates and eleven rows disappear. The 13 residuals in 26
variables have degree at most four, and the sum of squares has coefficient
one at $a^4f^4$. These are exact integer-polynomial zero-set bijections,
with no hidden division. Explicit restorations are contained in each JSON.

\paragraph{Alternative and checks.}
Source 69's weighted-cone guard also composes, giving the same three sizes.
Its canonical fixed-anchor bounds survive the same parity tuning, but its
six numerals require $881{,}696$ raw bytes, so the determinant is preferred.
The checker verifies all eight packets against source 70's literal packets,
all restoration identities and authentic coefficient payloads, liveness,
exact degrees, and altered-guard negative controls.
Finite replay is evidence for those algebraic and fixed-data claims, not a
substitute for the unbounded mathematical arguments or a new formal proof
of the repository's internal consistency biconditional.


\section{A complete compact equation in twenty-four integers}
\label{sec:compact}
This section prints a common base for the smaller projections, so that the
construction is more than a table of claimed sizes.
Take the six fixed integers $\beta,K_0,M_C,M_F,A_0,B_0$ bound in the preceding
section, with $A_0=2^{1113}$ and $B_0=2^{666104}$.
The 24 unrestricted unknowns are
\[
 J_3,F,\alpha,z,f,h,r,v,w_3,\tau,\eta,y,\gamma,o,j,E,c,a,V,I,t,p_1,p_2,p_3.
\]
The aliases $v,t,p_i$ stand for \texttt{v\_scale}, \texttt{pair\_root}, and
\texttt{pair\_pos}$i$ in the literal JSON. Define
\begin{align*}
 J&=3J_3+1,&q&=\beta J+1,&W&=4(3w_3+1),\\
 X&=qW,&Y&=(2v+1)q^3,&C&=q-F-\alpha,\\
 k&=1-V+jc+hE,&H&=4a+3,&\Delta&=a^2+H,\\
 R&=[(q-1)(q-F)+\alpha](q^2-1)+(M_C+qM_F)J,\\
 D_0&=(A_0F)(B_0C)-(q-A_0F-B_0C)^2,&T_0&=\eta(k-\eta),\\
 \Pi&=1+p_1^2+p_2^2+p_3(p_3+1).
\end{align*}
Every displayed symbol on a left-hand side is only an abbreviation.
The twelve residuals are
\begin{align*}
 e_1&=a-E-Y,\\
 e_2&=E(X+1)-aX,\\
 e_3&=c-k(a-E)-\eta,\\
 e_4&=\tau^2-E(c-\eta)\bigl[k+E(c-\eta)\bigr]-1,\\
 e_5&=(X+ac+\gamma H)^2-\Delta c^2-1,\\
 e_6&=jc-V-R,\\
 e_7&=(K_0+W)C-F-z(q-1),\\
 e_8&=V+c-fo,\\
 e_9&=I-(2r+1)c^2,\\
 e_{10}&=I^2-\Delta(f^2-1),\\
 e_{11}&=I^2V^2-(I^2-1)y^2-1,\\
 e_{12}&=D_0T_0-(t^2-D_0-T_0+1)^2-\Pi.
\end{align*}
Then
\[
 \boxed{P_{24}=e_1^2+e_2^2+\cdots+e_{12}^2=0,
 \qquad\boldsymbol{x}\in\Z^{24}.}
\]
It has exact degree 12. The new last row forces $D_0>0$ and $T_0>0$
by the universal simultaneous-positivity lemma below. It replaces two
separate ternary positivity witnesses by four witnesses, without changing
the other arithmetic coordinates.

For the 25-variable octic, add an unrestricted $\zeta$, add the residual
$\zeta-k+\eta$, and use $T_0=\eta\zeta$ only in the last row. Keep $k$ and
all other rows exactly as printed. This is an existential extension whose
projection is the 24-variable equation, but its degree is8 because the last
row is now quartic rather than sextic. The coefficient of $t^8$ is 1.
The independent root script \texttt{check\_displayed\_system.py} expands
and compares every printed row with both literal JSON artifacts.

The pair lemma and the further projections are proved next. In particular,
when a row is replaced using another retained equation, we assert equality
of zero sets, not equality of the entire sum-of-squares polynomials.
\section{Joint positivity and a projection family}
\label{sec:joint-projections}

All variables in this section range over unrestricted integers.  The six fixed
numerals are the unchanged authentic source-34 tuple and the two determinant
scales.  Write $d=U_cV_c-(q-U_c-V_c)^2$, where
$U_c=2^{1113}F$ and $V_c=2^{666104}(q-F-\alpha)$.
The separate symbol $W=4(3w_3+1)$ denotes the scale in the compact
arithmetic system.  These are polynomial abbreviations, not additional unknowns.

\subsection{An elementary universal two-sign lemma}

The following construction is the joint-positivity lemma of source 68,
\emph{Joint Positivity and Free Lucas}, Sections 1--2 of its
\texttt{PAIR\_GUARD\_PROOF.md}.  We claim its composition with the new
determinant/residue system, not the lemma as a new discovery.

Put
\[
 G(x,y,t)=xy-(t^2-x-y+1)^2,
 \qquad \operatorname{Pos}(p_1,p_2,p_3)
 =1+p_1^2+p_2^2+p_3(p_3+1).
\]
For integers $x,y$,
\[
 x>0\ \&\ y>0
 \quad\Longleftrightarrow\quad
 \exists t,p_1,p_2,p_3\in\mathbb Z:
 G(x,y,t)=\operatorname{Pos}(p_1,p_2,p_3).       \tag{JP}
\]
Indeed, $G>0$ implies $xy>0$.  If $x=-a,y=-b$ with $a,b\geq1$,
then $(t^2+a+b+1)^2\geq(a+b+1)^2>ab$, a contradiction.  This proves
soundness for every signed assignment, including arbitrary $t$.

Conversely, for $x,y\geq1$ set $N=x+y-1$ and $m=\lfloor\sqrt N\rfloor$.
Choose $t=m$ when $N-m^2\leq m$, and $t=m+1$ otherwise.  Then
$|t^2-N|\leq m$.  The error is zero when $N$ is a square; otherwise
$(t^2-N)^2\leq m^2<N\leq xy$, since $xy-N=(x-1)(y-1)\geq0$.
Thus $G>0$.  Every positive integer $n$ is a value of
$\operatorname{Pos}$: apply the three-square theorem to $4n-3$, which is
positive and congruent to $1$ or $5$ modulo $8$.  Exactly one square is odd,
so the representation has the form
$(2p_3+1)^2+(2p_1)^2+(2p_2)^2$.  Division by four gives the asserted
identity.  The floor and square root are only a witness-selection recipe
in this proof; neither appears in a polynomial.

\subsection{The 25-variable octic and its 24-variable projection}

The 26-variable determinant/residue octic has eleven arithmetic rows and
two private ternary positivity rows asserting
\[
 d>0,\qquad \eta\zeta_0>0,
 \qquad \zeta_0=1-V+jc+hE-\eta.
\]
Delete those six private coordinates and two rows.  Appending
\[
 d\eta\zeta_0-(t^2-d-\eta\zeta_0+1)^2
       -\operatorname{Pos}(p_1,p_2,p_3)=0             \tag{J24}
\]
gives 24 variables and 12 rows.  Its joint row has degree six.  Its top
part is $-\eta^2(hE+jc)^2$, so the sum of squares has exact degree twelve,
with coefficient $+1$ at $\eta^4h^4E^4$.

To keep degree eight, retain one extra coordinate $\zeta$, impose
$\zeta=\zeta_0$, and replace $\zeta_0$ by $\zeta$ only in the joint row.
Both $d$ and $\eta\zeta$ are now quadratic, so the new row has degree four.
This gives 25 variables and 13 rows.  Its coefficient at $t^4$ is $-1$;
no other row mentions $t$.  Hence the final polynomial has coefficient
$+1$ at $t^8$ and exact degree eight.  The eleven old arithmetic rows
are unchanged and may still contain $\zeta_0$; the restoring row makes
this intentional choice equivalent to using $\zeta$ everywhere.

For either construction, (JP) restores the two old positive quantities,
and their ternary representations restore all six deleted coordinates.
Conversely, every parent zero supplies those quantities and therefore the
four new witnesses by the explicit recipe above.  Thus this step preserves
existence at every fixed assignment of retained arithmetic coordinates.
It requires no additional canonical restriction.  Canonical completeness
is inherited only where it is actually used upstream, in the determinant
and residue constructions.

\subsection{A degree-sensitive row operation}

Use the compact system's abbreviations
\[
 X=4q(3w_3+1),\quad Y=(2v_{\rm scale}+1)q^3,
 \quad k=1-V+jc+hE,\quad\Delta=a^2+4a+3.
\]
The relevant rows are
\[
\begin{gathered}
 a-E-Y=0,\quad E-X(a-E)=0,\quad c-\eta-k(a-E)=0,\\
 V+c-fo=0,\quad I-(2r+1)c^2=0,\\
 S=I^2-\Delta(f^2-1)=0,\quad
 A=I^2(V^2-y^2)+y^2-1=0.
\end{gathered}
\]
Define $A_{\rm red}=\Delta(f^2-1)(V^2-y^2)+y^2-1$.  The exact identity
\[
       A-A_{\rm red}=(V^2-y^2)S                 \tag{RI}
\]
allows either auxiliary row while $S$ is retained.  It is a two-sided
zero-set row operation, not equality of the two free polynomials.

Every elimination below uses a monic defining row and an integer-polynomial
inverse.  Substitution maps old zeros to new zeros; restoring the deleted
coordinate makes its defining row zero and recovers every other old row.
There is no division, sign assumption, or new Pell existence theorem.

\begin{itemize}
\item At 23 variables, eliminate $V=fo-c$ and keep $A$.  The auxiliary and
joint rows have degree six, so the final degree is twelve.  Alternatively,
use (RI), eliminate $I=(2r+1)c^2$, and keep $V$; the degree is again twelve.
The latter coordinate choice has a smaller expanded support in the exact
independent audit, but no claim of circuit minimality follows.
\item At 22 variables, use (RI) and eliminate both $V,I$.  The reduced
auxiliary has degree eight and top term $a^2f^4o^2$, giving final degree
sixteen.
\item At 21 variables, instead retain the original $A$ and eliminate
$V,I,a$, with $a=E+Y$.  The main, strong, and auxiliary rows have degree at
most ten; all other rows have degree at most six.  This gives degree twenty.
\item At 20 variables, still keep the original $A$ and retain $\eta$.
The first two rows give $E=XY$ and $a=(X+1)Y$.  Together with the $V,I$
restorations, these leave eight rows.  The joint row has degree sixteen,
the main/strong and first Pell rows have degree at most fourteen, and the
original auxiliary has degree ten.  The final degree is thirty-two.
\item At 19 variables, also eliminate $\eta=c-Yk$.  The joint row has
degree forty-four and the other six rows have degree at most thirty-four.
The final degree is eighty-eight.
\end{itemize}

The choice of auxiliary row is essential for the degree ledger.
Keeping $A_{\rm red}$ at the 20-variable endpoint would needlessly raise
its degree.  These endpoints may all be defined directly from the compact
24-variable source by the stated substitutions and, only where specified,
the row identity (RI).

\subsection{Exact-degree certificates on the authentic slice}

For brevity in this table only, write $J=J_3$, $v=v_{\rm scale}$ and
$w=w_3$.  Each entry is an exact coefficient of the final sum of squares.
\begin{center}
\begin{tabular}{ccll}
Variables & Degree & Monomial & Coefficient\\\hline
25 & 8  & $t^8$ & $1$\\
24 & 12 & $\eta^4h^4E^4$ & $1$\\
23 & 12 & $\eta^4h^4E^4$ & $1$\\
22 & 16 & $a^4f^8o^4$ & $1$\\
21 & 20 & $f^4v^4J^{12}$ & $2^4 3^{12}\beta^{12}$\\
20 & 32 & $\eta^4h^4w^4v^4J^{16}$ & $2^{12}3^{20}\beta^{16}$\\
19 & 88 & $h^8w^8v^{16}J^{56}$ & $2^{32}3^{64}\beta^{56}$
\end{tabular}
\end{center}
All are positive because the authentic fixed $\beta=2^{1064921}-1$ is
positive.  The 23-variable certificate applies to both coordinate choices.

Here is an unbounded check of the last three entries.  The leading forms
are
\[
 X_*=36\beta Jw,\qquad Y_*=54\beta^3J^3v.
\]
At 21 variables the strong row has top term $-Y_*^2f^2$.
The main row has no $f$, and the original auxiliary has no $J$, so the
displayed coefficient can come only from the strong square.
At 20 variables, $E_*=X_*Y_*=1944\beta^4J^4wv$ and the joint top form is
$-\eta^2h^2E_*^2$; its square is the only contribution of degree 32.
At 19 variables,
\[
 \eta_*=-hX_*Y_*^2=-2^4 3^8\beta^7 hJ^7wv^2.
\]
Since $\deg\eta>\deg k$, the joint top form is $-\eta_*^4$.
Its square gives the displayed degree-88 monomial, with no other residual
of sufficiently high degree to contribute.  These arguments preclude
cancellation at the actual numeral values, rather than merely asserting
a generic gate-degree bound.

The 19-variable coordinate list is exactly
\[
 J_3,F,\alpha,z,f,h,r,v_{\rm scale},w_3,\tau,y,\gamma,o,j,c,
 t,p_1,p_2,p_3.
\]
In particular $X,Y,E,a,V,I,\eta$ and the coefficient ports are not hidden
quantified variables.

\subsection{Certificate scope}

All residuals are serialized as ordinary, completely inlined polynomial
expressions.  High-degree packets retain factorization to avoid a wasteful
expanded sum of squares.  A separate standard-library checker authenticates
the compact source identities, the row operation, every deleted defining
row, the literal surviving-row substitutions, the declared variable set,
and the exact homogeneous coefficient certificates.  It passes both
normal and optimized Python execution.  Finite tests are not substituted
for the universal proof of (JP), the inverse maps, or the inherited
compiler, Pell, and refutation-calibration theorems.

Consequently the ordinary mathematical existence equivalence with a finite
ZFC refutation extends to every listed packet.  No new kernel-checked
Lean/Rocq theorem or encoded object-ZFC derivation is claimed, and this
bounded family of variable/degree tradeoffs is not an optimality theorem.


\section{Literal shared circuits for the integer family}
The residual systems also admit small shared binary evaluators. The fixed
numeral leaves are free in the same counting convention as the input
packages; every live scalar multiplication, squaring, addition and subtraction
is charged. Registers are reused, not re-expanded into expression trees.
The following counts specify delivered circuits, not optimal circuits.
\begin{center}
\begin{tabular}{rrrr}
\toprule
Integer unknowns & Exact degree & One polynomial & Residual system\\
\midrule
25&8&127&102\\
23&12&121&100\\
22&16&118&99\\
21&20&115&98\\
20&32&112&97\\
19&88&109&96\\
\bottomrule
\end{tabular}
\end{center}
The system column charges the union of live gates needed to evaluate all
residuals, not the sum of separately recomputed expressions. The polynomial
column additionally forms and sums their squares. The 23-variable circuit
here is the variant eliminating $V$; the distinct variant eliminating $I$
has a smaller expanded support and is supplied separately.

The circuit proof is literal polynomial equality, stronger than equality
of zero sets. It validates each restored-coordinate definition and each
residual with named symbolic intermediate expressions, then structurally
inlines those expressions to match the corresponding JSON packet.
For example, the nineteen-variable first norm simplifies using the exact
small-ring identity
\[
 (XY)(kY)\bigl(k+(XY)(kY)\bigr)
 =XY^2(XY^2+1)k^2.
\]
Such identities avoid expanding the entire degree-88 sum of squares while
still proving an identity for every integer assignment. Random modular
evaluation alone would not suffice. The independent checker separately
checks topology, liveness, all declared inputs, counted binary operations,
and the symbolic residual bindings. No modular test is substituted for
an unbounded circuit identity.

The degree-88 equation in 19 unrestricted integers thus uses 109 operations,
with the same six numerical payloads as the other determinant constructions.
That improves the selected older 24-integer degree-1936/161-operation point
on variables, degree and this circuit-cost convention. Coefficient payloads
and proof dependencies are still separate comparison axes, so this is not
a global claim of the simplest possible equation.

\section{Denominator-free and centered three-root transport}
\label{sec:pell-three-root}

Two new eleven-positive-variable endpoints improve the inspected
degree-$325$, $113$-operation construction of source~70:
\[
 \boxed{\deg P_{11}=277,\quad 111\text{ operations}},\qquad
 \boxed{\deg \widetilde P_{11}=257,\quad 114\text{ operations}}.
\]
The first deletes a denominator guard using the adjacent-coefficient
Pell descent of source~69.  The second additionally centers the auxiliary
root and proves transport positivity by a new signed bootstrap.  They are
different degree/cost choices.  Every multiplication by a numeral and every
square is charged; fixed numeral leaves and register reuse are free.
Neither construction changes the known nine-positive-variable bound or
claims global optimality.  All native compiler and literal
machine/refutation calibrations remain named inherited dependencies.

\subsection{A denominator guard that can be deleted}

The following elementary fact replaces the prime-support argument in
source~70.

\paragraph{Adjacent-coefficient descent.}
If $Q>1$ and $V$ are integers, $y$ is rational, and
\[
 QV^2-(Q-1)y^2=1,
\]
then $y$ is an integer.  Indeed, replace both coordinates by their
absolute values and write $y=a/m$ in lowest terms.  Then $m^2\mid Q-1$.
For $V>1$ apply
\[
 V'=(2Q-1)V-2(Q-1)y,\qquad
 y'=(2Q-1)y-2QV.
\]
Here $V'$ is integral, the same norm is preserved, and $y>V$ and
$QV>(Q-1)y$ give $-V<V'<V$.  The norm excludes $V'=0$.
Since $2Q-1\equiv1\pmod m$, the reduced denominator of $y'$ is still
$m$.  Strong induction reaches $V=1$, where $y=1$, and proves $m=1$.

Thus, after a sum-of-three-roots argument has made the auxiliary root
rational, no denominator constraint involving the first transport modulus
is needed.  In particular source~70's $I=iq(q-1)c^2$ can be replaced by
$I=ic^2$, deleting two multiplication gates.  The unconditional separation
proof used only $I\ge c^2$ and remains valid.

\subsection{Complete polynomials}

The eleven unknowns are
\[
 j,F,\alpha,i,T,X,Y,n,\eta,\zeta,\gamma\in\mathbb Z_{>0}.
\]
The four fixed numeral ports are the complete compact44 tuple
$\beta,K_0,M_C,M_F$, with $\beta=2^{590069}-1$ and $M_F$ already
shifted.  They are not free parameters; their authenticated payloads and
whole-tuple binding accompany the circuits.  Every other symbol below is
a polynomial register, not another quantified variable.  Define
\begin{gather*}
 J=3j+1,\qquad q=\beta J+1,\qquad C=q-F-\alpha,\\
 E=XY,\qquad k=\eta+\zeta,\qquad c=kY+\eta,\\
 a=E+Y,\qquad H=4a+3,\qquad \Delta=a^2+H=(a+2)^2-1,\\
 \ell=EkY,\qquad I=ic^2,\qquad M=q^3(q-1),\\
 R=((q-1)(q-F)+\alpha)(q^2-1)+(M_C+qM_F)J,\\
 \widehat R=q^3(q+F+\alpha),\qquad B_t=(qK_0+X)C-qF,\\
 D_\tau=\ell(\ell+k)+1,\qquad D_f=1+\Delta I^2,\\
 Q=\Delta^2I^2,\qquad \delta=Q-1,\qquad V=D_f(cT-R)-c,\\
 r_0=k-R-1,\qquad S=4qY,\qquad U=qS,\\
 D_*=U(X+ac)+H(U\gamma+SX-r_0),\\
 A_1=D_*^2-U^2(1+\Delta c^2).
\end{gather*}

For the degree-$277$ endpoint choose
\[
 L=3D_f(cT+\widehat R),\quad h_*=L,\quad
 W=1+(q-1)h_*,\quad
 Z=M(\ell+k+n+L)+(q-1)Y-q^3WB_t.
\]
For the degree-$257$ endpoint instead choose
\begin{align*}
 B_0&=cT+\widehat R,&L&=D_f+B_0,\\
 h_*&=4qD_f+L,&W&=1+(q-1)h_*,\\
 Z&=M(\ell+k+n+L+V)+(q-1)Y-q^3WB_t.
\end{align*}
The shared division-free final norm is
\begin{align*}
 A&=M^2D_\tau,&B_f&=M^2D_f,&C_y&=M^2(QV^2-1),\\
 T_N&=\delta^2AB_f,&R_N&=\delta C_y,\\
 J_N&=\delta(Z^2-A-B_f)-C_y,\\
 N_c&=J_N^2-4\bigl(T_N+R_N(A+B_f)\bigr),\\
 N_8&=N_c^2-64T_NR_NZ^2,\\
 P&=N_8+A_1+N_8A_1.
\end{align*}
The literal graphs are \texttt{p11\_positive\_degree277.json} and
\texttt{p11\_positive\_degree257.json}.

\subsection{Three-root separation and the degree-277 lift}

In proof only put $D_y=(QV^2-1)/(Q-1)$, and let
$t=\sqrt{D_\tau}$, $f=\sqrt{D_f}$ and $y=\sqrt{D_y}$ be positive.
For every positive input, independently of any equation or sign of $C$,
\[
 1<D_\tau<D_f<Q,\qquad |V|>1,\qquad D_y>V^2.
\]
The radicand $D_y$ cannot equal either integer $D\in\{D_\tau,D_f\}$:
that would give $Q(D-V^2)=D-1$ with $0<D-1<Q$.
There is also strict root dominance.  Write $b=cT-R$.
If $b=0$, then $V=-c$, $y<c+1$ and $f>t+y$.
If $b\ne0$, then $|V|\ge D_f-c>2f>f+t$, hence $y>f+t$.
These estimates use $I\ge c^2$, $c\ge3$ and
$f>(a+1)I\ge3c^2$.  Thus no signed sum of the three roots is zero.

The universal octic sum norm is
\[
 F_8(x;a,b,c)=
 \bigl[(x^2-a-b-c)^2-4(ab+ac+bc)\bigr]^2-64abcx^2.
\]
It is the product of the eight signed-root factors.  The displayed
$N_8$ is exactly $\delta^4F_8(Z;A,B_f,C_y/\delta)$.
At $N_8=0$, therefore,
\[
 Z/M=\tau_*+f_*+y_*,\qquad
 \tau_*^2=D_\tau,\quad f_*^2=D_f,\quad y_*^2=D_y.
\]
A rational signed sum of three distinct square roots, with no zero
signed sum, forces each root rational.  One exact extraction is as follows:
for signed roots $x_1,x_2,x_3$ with rational nonzero sum $s$ and squares
$a,b,c$, put
\[
 u=(s^2-a-b-c)/2,\quad
 v=(u^2-ab-ac-bc)/(2s),\quad
 x_1=\frac{a(a+u)}{as+v}.
\]
The denominator is $x_1(x_1+x_2)(x_1+x_3)\ne0$; cycle the formula.
Consequently $\tau_*,f_*$ are integers by rational-square integrality,
and $y_*$ is an integer by the preceding Pell descent.

For the $277$-degree construction, $L>f+y$ unconditionally.  Recover
\[
 r=n+L-f_*-y_*>0,\quad f_{\rm old}=|f_*|,\quad
 i_{\rm old}=i,\quad u=|y_*|-V>0.
\]
These are precisely the weighted source equations.  Conversely choose
canonical strong and auxiliary source witnesses first, then $L,W$ and
the transport witness.  The inherited weighted theorem gives
$r>qW-2k>L$, so $n=r-L+f+y>0$.  No stronger divisibility than
$c^2\mid I$ is required.

\subsection{The new centered sign bootstrap}

The centered endpoint cannot assert $r>0$ before proving $V>0$.
The following order resolves that issue.

First recover all three integral roots as above.  Since $M\mid Z$,
reduction modulo $q^3$ and $q-1$ gives
\[
 q^3\mid Y,\qquad q-1\mid B_t.
\]
No transport positivity is used here.  Also $A_1=0$ makes $D_*/U$
integral; classify its \emph{absolute value} as a positive Pell root.
Together with the first Pell root, this gives
\[
 c=\psi_{a+2}(p),\quad k=2\psi_{2XY^2+1}(m),\quad p,m\ge1.
\]
Here $c\ge3$ implies $p\ge2$.  The case $m=1$ would give $k=2$ and
$c=2Y+1<2(a+2)$, a contradiction.  Hence
\[
 c>8XY^3\ge8q^9. \tag{*}
\]

The authentic mask bounds give $|R|<\widehat R$, hence
$|V|<D_fB_0$, while $\delta>D_f$ and $D_f>f$.  Therefore
\[
 0<y-|V|<|V|/\delta<B_0.
\]
In particular $L+V>f+y$ if $V>0$.

Suppose now that $C\le0$.  Set $d=-C\ge0$ and
$B=-B_t=(qK_0+X)d+qF>0$.  Since $K_0\ge1$ and
$(M_C+qM_F)J<2q^2$, direct expansion gives
\[
 qB\ge q^2d+q^2F,\qquad R<qB+q^4+q^2.
\]
Writing $s=Y/q^3>0$, the root equality becomes
\[
 \tau_*+f_*+y_*=
 \ell+k+n+L+V+s+\frac{WB}{q-1}. \tag{**}
\]
If $V>0$, the right side is greater than $t+f+y$, impossible.
If $V<0$, then the integer $cT-R$ is nonpositive, so
$R\ge cT\ge c>8q^9$.  Thus $q^4+q^2<R/2$ and $R<2qB$.
Moreover $|V|<D_fR$, so $L+V-f-y>-2D_fR$, whereas
\[
 \frac{WB}{q-1}>4qD_fB>2D_fR.
\]
Again the right side of $(**)$ exceeds $t+f+y$.
The impossible case $V=0$ was already excluded.  This proves $C>0$.

Only now use the authentic bound $3q<R<q^4-2$.
Together with $(*)$ it gives $V>1$, and therefore
\[
 r=n+L+V-f_*-y_*\ge n+L+V-f-y>n>0.
\]
The old positive strong and auxiliary coordinates are recovered exactly
as before.  The weighted carrier soundness theorem now applies.

For completeness choose canonical strong and auxiliary data before
computing $L,h_*,W$.  Since $h_*>L+k$, the canonical transport variable
satisfies $r>qW-2k>L$.  Set
\[
 n=r-L-V+f+y=r-L+f+u>0.
\]
This gives the displayed centered norm.  The old positive main-carrier
coordinate $\gamma$ is unchanged.  This is an existence equivalence,
not a bijection between every old and new zero set.

Both $N_8$ and $A_1$ are squares modulo four.  Thus
$(1+N_8)(1+A_1)=1$ forces both factors to be $1$, since $-1$ would
require a residual $-2$.  The single equation is therefore equivalent
to the two residual equations.  Under the inherited interfaces, positive
solvability is equivalent to a finite ZFC refutation; negation gives
the stated consistency equivalence.

\subsection{Exact degrees and checks}

For the uncentered construction the exact degrees of
$I,D_f,Q,V,L,W,Z,N_c,N_8,A_1$ are respectively
\[
 5,14,18,18,18,19,24,132,264,13.
\]
The coefficient of
$j^{60}i^{24}\eta^{49}X^{42}Y^{85}\alpha^{16}\gamma$
in the degree-$277$ polynomial is
$2^7 3^8(3\beta)^{60}>0$.

For the centered construction put $z=Z-MV$ in proof only.  The exact
cancellation identity is
\[
 J_N=\delta\{z(z+2MV)-A-B_f\}-M^2(V^2-1).
\]
Here $\deg z=21$, $\deg Z=22$, and the five displayed summand degrees
are $60,61,34,40,44$.  Thus $\deg J_N=61$,
$\deg N_c=122$ and $\deg N_8=244$; the subtracted norm correction
and tail have degrees $102$ and $198$, respectively.  The coefficient of
\[
 j^{52}i^{24}\eta^{49}X^{38}Y^{85}F^4\alpha^4\gamma
\]
is $2^{19}(3\beta)^{52}>0$, proving exact degree $257$ after actual
fixed-port specialization.

Both standard-library builders reproduce the literal graphs byte for
byte.  Both exact checkers pass normally and under \texttt{python -O}:
prefix identities, the eight-factor norm, rational extraction identities,
main-carrier cancellation, centered cancellation, degree heads, selected
coefficients, gate ledgers and fixed coefficient hashes.  Independent
review receipts are kept separately.  These finite checks do not prove
the unbounded native/compiler interfaces or the consistency of ZFC.

\paragraph{Unfinished four-root direction.}
A direct ten-variable extension is not established.  Distinct radicands
and a nonzero full signed sum no longer ensure rationality root by root:
$\sqrt2+\sqrt8-\sqrt{18}+100=100$, with four distinct radicands and
no vanishing full signed sum.  Even a stronger separation argument would
not remove the main-residue versus transport-modulus obstruction from
the earlier attempted merger.  The negative controls in this component
record both issues; no ten-variable endpoint is claimed.

\section{Smaller local graphs for the direct checker}
\label{sec:direct-local-wave11}

The latest pinned repository supplies a complete numeric checker for the
original sealed-ZFC calculus, independent of K215 and Jones finalization.
Its $83{,}281$ raw and $76{,}313$ optimized primitive-recursive program nodes
are not polynomial variable counts.  This section improves concrete local
polynomial graphs used by that route, without changing its numeric codes.
Every visible variable and witness below ranges over $\mathbb N$, including
zero; every subtraction is ordinary integer-polynomial subtraction.

\begin{center}
\begin{tabular}{lrrr}
\hline
Local graph & Old witnesses/ops & New witnesses/ops & Degree\\
\hline
Max-square pair & $3/46$ & $2/39$ & $4$\\
Pair, strict squared rows & $3/46$ & $2/42$ & $4$\\
Either unpair projection & $4/46$ & $3/39$ & $4$\\
Predecessor & $0/16$ & $0/10$ & $4$\\
Tree left/right, list tail & $5/62$ & $4/49$ & $4$\\
If-zero conditional & $2/32$ & $1/16$ & $4$\\
Zero test & $1/20$ & $1/14$ & $4$\\
Actual list head & $7/76$ & $5/59$ & $4$\\
Pair, selector eliminated & $3/46$ (quartic) & $1/53$ & $6$\\
\hline
\end{tabular}
\end{center}
Operations count binary additions, subtractions and multiplications in the
literal expression tree, including the final addition of zero, as in the
baseline.  Repeated subtrees are counted repeatedly.  The artifacts also
provide separate shared-DAG counts and full expanded integer coefficients.
No local saving is advertised as a completed smaller consistency polynomial.

\subsection{Remove the maximum-component witness}
For natural $z,a,b,t,s$, put
\[
 B=t(t-1),\qquad L=a-b-s+t(2s+1),\qquad
 E=4(z-a-b)-(a+b+s)^2+t,
\]
\[
 \operatorname{Pair}(z,a,b;t,s)=B+L^2+E^2.
\]
Then
\[
 \exists t,s\in\mathbb N:\operatorname{Pair}=0
 \quad\Longleftrightarrow\quad
 z=\begin{cases}b^2+a,&a<b,\\a^2+a+b,&b\le a.\end{cases}
\]
Indeed $t(t-1)\ge0$ for every integer $t$, so a zero forces $t=0$ or
$t=1$ and $L=E=0$.  At $t=0$, $a=b+s$ and $a+b+s=2a$; hence
$z=a^2+a+b$.  At $t=1$, $b=a+s+1$ and $a+b+s=2b-1$; hence
$z=b^2+a$.  Conversely these two branches choose respectively
$(t,s)=(0,a-b)$ and $(1,b-a-1)$.  This covers ties and zero.
The literal quartic has two witnesses, $39$ tree operations and $52$
monomials.  Replacing $B$ by $B^2$ gives a strict sum of three squared
quadratic rows, with $42$ operations and the same zero set.  This latter
version fits the existing squared-row interface without a nonnegative-term
extension.  Hiding the missing component yields either unpair projection
with three witnesses.

Solving the last row for the selector also gives the one-witness sextic
\[
 T=(a+b+s)^2-4(z-a-b),\qquad
 T(T-1)+\bigl(a-b-s+T(2s+1)\bigr)^2=0.
\]
Here $T$ is an integer polynomial, not a new coordinate.  Its nonnegative
first summand forces $T\in\{0,1\}$; the preceding branch proof then applies.
This has $53$ tree operations and $86$ monomials, but raises degree to six.

\subsection{Share the predecessor and remove branch selectors}
The exact predecessor graph factors as
\[
 [n(p+1-n)]^2+[p(p+1-n)]^2=(n-p-1)^2(n^2+p^2).
\]
Its zero set is precisely $p=\max(n-1,0)$: either $n=p+1$ or $n=p=0$.
Adding this product to $\operatorname{Pair}(p,a,b;t,s)$ gives either tree
projection with four hidden coordinates and $49$ operations.  At $n=0$
both components are zero, so the right projection is also the actual guarded
list tail.

A separate selector is unnecessary for a conditional.  The quartic
\[
 [n(c-b)]^2+[(c-a)(n-s-1)]^2=0
\]
with one natural witness $s$ is equivalent to $c=a$ if $n=0$, and $c=b$
otherwise.  At zero, $n-s-1=-s-1\ne0$; at positive $n$, the first row
forces $c=b$, and completeness chooses $s=n-1$.  Taking $a=1,b=0$ gives
the $14$-operation zero test.  If an exact predecessor $p$ is already
available, set $s=p$ and allocate no extra slack.

For the actual found-head convention, with visible $n,c$ and witnesses
$p,a,b,t,s$, use
\[
 \operatorname{Head}=\operatorname{Pair}(p,a,b;t,s)
 +(n-p-1)^2(n^2+p^2+c^2)+[n(c-a-1)]^2.
\]
At $n=0$, its middle term forces $p=c=0$ and the pairing term forces
$a=b=0$.  At $n>0$, it forces $p=n-1$, $c=a+1$, and the correct unpair
components.  Conversely these choices make every term zero.  Thus nil is
zero and a found zero-valued entry is one.  The ordinary quartic has five
witnesses, $59$ operations and $72$ monomials.

\subsection{Certificate status}
The package supplies literal trees and full coefficient lists for eleven
local polynomials.  Exact symbolic identities, counts and $1{,}845{,}448$
finite controls pass in ordinary and optimized Python modes; checks remain
active under \texttt{-O}.  Universal correctness follows from the arguments
above, not from the finite controls.  These new proofs have not been checked
by Lean, and the whole numeric checker, its bounded loops and its internal
nonstandard-code interpretation still require integration.  Natural witness
domains are essential; unrestricted integer witnesses need their own
naturalness encoding.

\section{Index-free inert-prime closure with a reused product scale}
\label{sec:index-free-crt-wave11}

Fix $m\ge1$, natural $n$, external natural parameters $x$, and fixed
nonnegative-coefficient polynomials $e(i,w;x),f(i,w;x)$.  Set $P=e-f$ and
choose a fixed $d\ge0$ bounding its degree in $i$; for $P=0$, $d=0$ is
allowed.  For a natural witness bound $U$, define
\[
\begin{split}
 H&=e(n,U,\ldots,U;x)+f(n,U,\ldots,U;x)+n+U+1,\\
 T&=H!,\qquad D=4T,\qquad S=D^n n!,\qquad M=DS+1,\\
 B&=\bigl(S\tbinom{S+n}{n}\bigr)\bmod M.
\end{split}
\]
Writing $P(i,w;x)=\sum_{k=0}^dP_k(w;x)i^k$, set
\[
 R_D(A;x)=\sum_{k=0}^dP_k(A;x)D^{d-k}(1-D)^k.
\]
This is an ordinary polynomial.  Its inverse-looking shorthand
$D^dP(D^{-1}-1,A;x)$ introduces no actual inverse or variable exponent.
For a fixed nested binary norm $N$, obtained by repeatedly replacing a pair
of subexpressions by the sum of their squares, the exact certificate is
\[
 \exists U,A_1,\ldots,A_m\in\mathbb N:\quad
 B\mid N\left(\tbinom{A_1}{U+1},\ldots,
                 \tbinom{A_m}{U+1},R_D(A;x)\right).
\]
It is equivalent to $\forall i<n\ \exists w\in\mathbb N^m:P(i,w;x)=0$.
Binomial is total, with $\binom Ak=0$ for $A<k$.  Neither factorials nor
binomials nor $D^n$ are being called ordinary polynomial notation.

\paragraph{Why the product needs only one variable-power graph.}
For $m_i=D(i+1)-1$,
\[
 S\binom{S+n}{n}=\prod_{j=1}^n(DS+Dj)
 \equiv\prod_{j=1}^n(Dj-1)\pmod{DS+1}.
\]
For $n>0$ the positive product is smaller than $D^n n!=S<M$; at $n=0$
it is one, still below $M$.  Hence $B=\prod_{i<n}m_i$.  The former
independent bound $(Dn)^n+1$ is unnecessary: reuse $S$ itself.  This removes
one complete variable-power graph, not merely its displayed abbreviation.

\paragraph{Why the encoded index is unnecessary.}
The integer polynomial $R_D-D^dP(i,A;x)$ is divisible by $D(i+1)-1$,
by the difference-of-powers identity.  Thus modulo $m_i$, the cleared
residual is exactly $D^dP(i,A;x)$.  The modulus already encodes the index
as $D^{-1}-1$, so no CRT coordinate $c$ or index norm leaf is needed.

Each $m_i$ is $3$ modulo $4$, all its prime factors exceed $H$, and the
$m_i$ are pairwise coprime: a common divisor divides $j-i$, hence $H!$,
but is coprime to $D$.  Choose a prime $p_i\equiv3\pmod4$ dividing $m_i$.
Anisotropy of a binary sum of squares over $\mathbb F_{p_i}$ makes every
norm leaf zero.  The binomial leaves recover residues $w_j\in[0,U]$;
$D$ is invertible, so the last leaf gives $P(i,w;x)=0\pmod{p_i}$.
Its original absolute value is smaller than $H<p_i$, hence it is zero
as an integer.  The bound is applied after cancelling $D^d$, not to the
large cleared residual.  Conversely CRT representatives of actual witnesses
make every leaf divisible by every $m_i$, proving completeness.  The empty
bound $n=0$ is handled by $B=1$.

\paragraph{An exposed, conditional witness ledger.}
Suppose chosen factorial, power and total-binomial graph polynomials use
respectively $\nu_F,\nu_P,\nu_B$ internal witnesses beyond visible arguments
and output.  Expose graph outputs
\[
 T=H!,\quad F=n!,\quad E=D^n,\quad C=\binom{EF+n}{n},\quad
 L_j=\binom{A_j}{U+1}.
\]
Use ordinary expressions $S=EF$, $M=DS+1$, and natural $B,K,r,V$ satisfying
\[
 SC=B+MK,\qquad M=B+r+1,\qquad N(L_1,\ldots,L_m,R_D)=BV.
\]
There are exactly $2m+9$ exposed coordinates.  The complete conditional
count is
\[
 2m+9+2\nu_F+\nu_P+(m+1)\nu_B,
\]
versus $2m+11+2\nu_F+2\nu_P+(m+1)\nu_B$ for the preceding interface
under the same convention.  The saving is exactly $\nu_P+2$, but no
numerical values for these uncompiled graph costs have been asserted.
This is a proved compiler simplification, not yet a smaller closed ordinary
consistency polynomial.  Further quartic flattening and internal formalization
have separate costs and proof obligations.


\section{Comparison and what has not been achieved}
The completed native family is a synthesis of previously separate tools,
not a claim that each ingredient was invented here. The new mathematical
obligations are their simultaneous canonical completeness, the ordering of
sign and rank recovery, and the exact inverse reconstruction after each
coordinate deletion. A superficially similar substitution can fail if it
uses a canonical bound on arbitrary signed inputs, or if it forgets an
integrality condition. The proof and independent checkers address these
specific risks.

The lower-variable integer projections keep the same six fixed payloads.
They do not claim reduced maximum expanded coefficient height or a minimum
arithmetic circuit. Their definitions are factored where expansion is
unhelpful; those registers still introduce no unknowns.
Historical integer points include 25 variables at degree 920,26 at degree 214,
and 24 at degree 1936. The new low-degree integer tradeoffs should be compared
with those domains, not with the existing25-positive sextic or20-positive
octic. Ten unrestricted variables and nine positive variables already exist
at much larger degrees, so our variable counts are not absolute records.

The 11-positive degree277 construction improves both degree and gate count
relative to the inspected325/113 predecessor. Its centered companion is a
separate degree/cost tradeoff. Attempts to merge a fourth Pell root and reach
ten low-degree positive unknowns still encounter preservation of the main
root and transport residue channels. Radical rationality alone would not
repair that obstruction. No such ten-variable theorem is claimed.

The direct compiler optimizations preserve the existing max-square pairing
and zero-test semantics. A different Cantor coding could make pairing itself
quadratic with no private auxiliary coordinate, but changing a whole proof
checker requires a recursive code translation and a correctness theorem.
That possibility is not a drop-in replacement for the current frozen
checker. The CRT interface similarly still requires complete polynomial
graphs for its remaining special functions. Neither local node counts nor
uncompiled graphs are counted as final consistency-polynomial sizes.

\section{Reproducibility and the formal proof boundary}
The release contains editable TeX, compiled PDF, literal/factored equation
specifications, complete coefficient bindings, source pins and inherited
proof texts, normal and optimized finite checks, and independent audits.
Original input ZIPs are preserved for authentication and replay.
Run \texttt{python3 verify.py} in an extracted release; add
\texttt{--replay-originals} to rerun the original incoming verifiers as well.
Run \texttt{python3 negative\_controls.py} for the root integrity controls.
The replay works in an isolated temporary copy and compares selected rebuilt
artifacts byte for byte. It does not decide whether any consistency equation
has a solution.

There are three separate assurances. The exact algebra checks certify the
finite displayed transformations and literal artifacts. The accompanying
unbounded arguments establish soundness and completeness of the new steps.
The inherited compiler/Pell/refutation interfaces provide the semantic
calibration to finite ZFC proofs. This packet does not newly formalize their
entire composition or export an encoded object-language ZFC biconditional.
The current upstream formal developments are reported with their own scope;
no local diagnostic is used to upgrade an upstream claim to a fresh kernel
check.

\paragraph{Source references.}
The project is Vladimir Reshetnikov's ProveIt repository,
\url{https://github.com/VladimirReshetnikov/ProveIt/tree/main/Computability/DiophantineConsistency}.
Full immutable hashes, archive Git blob IDs and SHA256 values are recorded
in the pin manifests. The reviewed comparison is
\url{https://github.com/VladimirReshetnikov/ProveIt/blob/2605669a532342f153b13eb6b4a0af4b54f5fa37/Computability/DiophantineConsistency/PROPOSAL_REVIEW_WAVE10.md}.
Authored proofs identify the relevant original archive and theorem at each
import boundary. All copied sources retain their notices and provenance.
\end{document}
